\documentclass[11pt]{article}
\usepackage[margin=2.3cm]{geometry}
\usepackage{amsmath,amssymb}
\usepackage{booktabs}
\usepackage{array}
\usepackage{enumitem}
\usepackage[hidelinks]{hyperref}
\usepackage{lmodern}
\setlength{\parskip}{4pt}
\setlength{\parindent}{0pt}

\newcommand{\ind}[1]{\mathbf{1}\!\left[#1\right]}
\DeclareMathOperator*{\argmax}{arg\,max}

\title{The PEARL Machine: the Full Math, Layer by Layer}
\author{Repository \texttt{Machine-}, branch \texttt{ccr-063c8579-mwbwim}}
\date{Data: BigONE PEARL-USDT, 25 Jun 2026 -- 20 Nov 2026}

\begin{document}
\maketitle

\section*{The stack at a glance}

The program is ten layers. Each layer only uses the output of the layer above it.

\begin{center}
\renewcommand{\arraystretch}{1.25}
\begin{tabular}{@{}r l l l@{}}
\toprule
 & Layer & Input & Output \\
\midrule
1 & Data & BigONE minute candles & prices $o_t,h_t,l_t,c_t$ \\
2 & Costs & fee, buy/sell gap & $\kappa$ per side \\
3 & Features & prices & $U_t,\ D_t,\ R_t$ \\
4 & Rule & features, setting $(b,s)$ & buy signal $\beta_t$, sell signal $\sigma_t$ \\
5 & Simulator & signals & money $V_E(b,s)$ \\
6 & Landscape & 40{,}000 settings & matrix $\mathbf{V}_E$ \\
7 & Smoothing & $\mathbf{V}_E$ & $\widetilde{\mathbf{V}}_E = K\,\mathbf{V}_E K^{\top}$ \\
8 & Pick & $\widetilde{\mathbf{V}}_E$ & setting $\theta^*(E)$ \\
9 & Machine & picks every $\Delta$ minutes & final money $F(\Delta)$ \\
10 & Interval sweep & $\Delta=1,\dots,1440$ & safest $\Delta$ \\
\bottomrule
\end{tabular}
\end{center}

\section{Data}

Time $t$ counts minutes. The data are BigONE PEARL-USDT one-minute candles. For each
minute: open $o_t$, high $h_t$, low $l_t$, close $c_t$.
The first 240 minutes are used only to start the 4-hour features, so the machine runs on
\[
t = 0,1,\dots,T-1,\qquad T = 212{,}932 \text{ minutes (25 Jun 20:38 to 20 Nov 17:29 UTC)}.
\]
First close $c_0 = 0.68$, last close $c_{T-1} = 1.9515$.

\section{Trading cost}

Each buy or sell pays the exchange fee $f$ plus half the gap between the buy price and the
sell price, $g$:
\[
f = 0.2\%,\qquad g = 0.5\%,\qquad \kappa = f + g = 0.7\% \text{ per side}.
\]
A buy at price $p$ costs $p(1+\kappa)$ per coin; a sell at price $p$ returns $p(1-\kappa)$ per coin.
One buy-then-sell at the same price loses
\[
1 - \frac{1-\kappa}{1+\kappa} \approx 2\kappa = 1.4\%.
\]

\section{Features (what the rule looks at)}

With a 4-hour window of $240$ minutes:
\begin{align*}
L_t &= \min_{0\le j<240} l_{t-j} && \text{(4-hour low)}\\
H_t &= \max_{0\le j<240} h_{t-j} && \text{(4-hour high)}\\
U_t &= 100\left(\frac{c_t}{L_t}-1\right) && \text{(\% above the 4-hour low)}\\
D_t &= 100\left(1-\frac{c_t}{H_t}\right) && \text{(\% below the 4-hour high)}\\
R_t &= 100\left(\frac{c_t}{c_{t-240}}-1\right) && \text{(\% change over 4 hours)}
\end{align*}

\section{The rule: buy the crash, sell the pump}

A setting is a pair $\theta=(b,s)$: the buy \% and the sell \%. Two constants are fixed:
$a = 0.23\%$ (``at the low'') and $n = 1.8\%$ (``near the high'').
\begin{align*}
\beta_t(b) &= \ind{U_t \le a}\cdot \ind{R_t \le -b} && \text{buy: fell $b$\% in 4 h and sits at the 4-hour low}\\
\sigma_t(s) &= \ind{D_t \le n}\cdot \ind{U_t \ge s} && \text{sell: $s$\% above the 4-hour low and near the 4-hour high}
\end{align*}

\section{The simulator (one setting)}

State at minute $t$: dollars $x_t$ and coins $q_t$. Start: $x_0 = 1000$, $q_0 = 0$.
A decision made at minute $t$ fills at the next minute's close $c_{t+1}$, so the rule never
sees the future. For $t = 0,\dots,E-2$:
\[
(x_{t+1},q_{t+1}) =
\begin{cases}
\left(0,\ \dfrac{x_t}{c_{t+1}(1+\kappa)}\right) & \text{if } q_t = 0 \text{ and } \beta_t(b)=1 \quad\text{(buy)}\\[2ex]
\left(q_t\,c_{t+1}(1-\kappa),\ 0\right) & \text{if } q_t > 0 \text{ and } \sigma_t(s)=1 \quad\text{(sell)}\\[1ex]
(x_t,\ q_t) & \text{otherwise.}
\end{cases}
\]
The money at minute $E$ counts coins as if sold then:
\[
V_E(b,s) = x_{E-1} + q_{E-1}\,c_{E-1}(1-\kappa).
\]

\paragraph{Same thing as a product.} Let the buy minutes be $\tau_1<\tau_2<\dots$ and the sell
minutes $\varsigma_1<\varsigma_2<\dots$, found by jumping:
\[
\tau_1 = \min\{t:\beta_t=1\},\qquad
\varsigma_k = \min\{t>\tau_k:\sigma_t=1\},\qquad
\tau_{k+1} = \min\{t>\varsigma_k:\beta_t=1\}.
\]
Then, for $N$ finished trades,
\[
V_E(b,s) = 1000\prod_{k=1}^{N}\frac{c_{\varsigma_k+1}\,(1-\kappa)}{c_{\tau_k+1}\,(1+\kappa)}
\quad\times\quad
\begin{cases}\dfrac{c_{E-1}(1-\kappa)}{c_{\tau_{N+1}+1}(1+\kappa)} & \text{if still holding}\\[1.5ex] 1 & \text{otherwise.}\end{cases}
\]
Each trade multiplies the money by (sell price $\div$ buy price) $\times\ (1-\kappa)/(1+\kappa)$.
This is why many small trades lose: each one pays about $1.4\%$.

\section{The landscape (40,000 settings at once)}

Both the buy \% and the sell \% take values on the grid
\[
\mathcal{G} = \{0.5,\ 1.0,\ 1.5,\ \dots,\ 100\},\qquad |\mathcal{G}| = 200,
\]
so there are $200\times200 = 40{,}000$ settings. The landscape at minute $E$ is the matrix
\[
\mathbf{V}_E \in \mathbb{R}^{200\times 200},\qquad
(\mathbf{V}_E)_{ij} = V_E(b_j,\ s_i),\qquad b_j = s_j = 0.5\,j .
\]
Rows are sell \%, columns are buy \%. A column whose buy \% was never reached stays at $1000$
(it never bought).

\section{Smoothing}

A Gaussian blur about $1\%$ wide (two grid steps), clipped at the edges. Weights:
\[
w_k = \frac{e^{-k^2/(2\cdot 2^2)}}{\sum_{m=-6}^{6} e^{-m^2/(2\cdot 2^2)}},\qquad k=-6,\dots,6 .
\]
Blur matrix $K\in\mathbb{R}^{200\times200}$:
\[
K_{ij} = \sum_{k=-6}^{6} w_k\ \ind{\operatorname{clip}(i+k,\,1,\,200) = j}.
\]
Blurring rows and columns at once:
\[
\widetilde{\mathbf{V}}_E = K\,\mathbf{V}_E\,K^{\top}.
\]
Each smoothed value is a weighted average of the settings within about $1\%$ around it. A
wide hill survives; a thin lucky spike gets flattened.

\section{The pick}

\[
(i^*,j^*) = \argmax_{i,j}\ (\widetilde{\mathbf{V}}_E)_{ij},
\qquad
\theta^*(E) =
\begin{cases}
(b_{j^*},\ s_{i^*}) & \text{if } (\widetilde{\mathbf{V}}_E)_{i^*j^*} > 1000\\
\varnothing & \text{otherwise (stay in USDC).}
\end{cases}
\]
$\theta^*(E)$ uses only minutes before $E$. On day one almost nothing has traded, so the
machine stays in USDC until some setting has made money.

\section{\texorpdfstring{The machine (re-pick every $\Delta$ minutes)}{The machine (re-pick every Delta minutes)}}

The machine holds the latest pick:
\[
\theta_t =
\begin{cases}
\varnothing & t < \Delta\\
\theta^*\!\left(\Delta\left\lfloor t/\Delta \right\rfloor\right) & t \ge \Delta.
\end{cases}
\]
It then runs the same simulator as Section 5, but with a setting that can change over time:
$\beta_t(b_t)$ and $\sigma_t(s_t)$ where $\theta_t=(b_t,s_t)$, and no trades while
$\theta_t=\varnothing$. Its final money is
\[
F(\Delta) = x_{T-1} + q_{T-1}\,c_{T-1}(1-\kappa).
\]

\section{The interval sweep and the safe choice}

All re-pick intervals $\Delta = 1,2,\dots,1440$ minutes are run, each starting with \$1{,}000 on
25 Jun with no data. To judge safety, look at the neighbours $\Delta-w,\dots,\Delta+w$:
\[
\text{worst}_w(\Delta) = \min_{|k|\le w} F(\Delta+k),\qquad
\text{spread}_w(\Delta) = \max_{|k|\le w} F(\Delta+k) - \min_{|k|\le w} F(\Delta+k).
\]
The safe choice is a $\Delta$ with a high worst case and a small spread. Every
$\Delta=20,\dots,39$ gives exactly the same result, so $\Delta^*=30$ sits in the middle:
\[
F(30) = \$4{,}733.99,\qquad \text{worst}_{10}(30) = \$4{,}733.99,\qquad \text{spread}_{10}(30) = \$179 .
\]

\section{The benchmark: just holding}

Buy at the first minute, sell at the last, paying the cost twice:
\[
\text{Hold} = 1000\cdot\frac{c_{T-1}}{c_0}\,(1-\kappa)^2 = 1000\cdot\frac{1.9515}{0.68}\cdot 0.993^2 = \$2{,}829.82 .
\]

\section{How it runs fast}

The trades of one setting up to minute $E$ do not depend on $E$ (the simulator never looks
ahead). So:
\begin{enumerate}[leftmargin=*]
\item Run all 40,000 settings over the whole history once and record every buy and sell as an
event $(t_e,\ \text{setting},\ \text{buy/sell},\ \text{amount})$: $191{,}132$ events.
\item Keep three arrays over all settings: dollars $\mathbf{x}$, coins $\mathbf{q}$, holding flag $\mathbf{h}$.
Walk forward in time; before minute $E$, apply every event with $t_e \le E-2$. Then
\[
\mathbf{V}_E = \mathbf{h}\odot \mathbf{q}\,c_{E-1}(1-\kappa) + (1-\mathbf{h})\odot\mathbf{x}.
\]
\item Compute $\theta^*(E)$ for every minute $E=1,\dots,T-1$ once: $212{,}931$ picks, each reading
all $40{,}000$ settings, about $8.5\times10^{9}$ values in total.
\item The pick at minute $E$ is the same for every machine, so all 1,440 machines reuse it.
Together they consult $\sum_{\Delta=1}^{1440}\lfloor T/\Delta\rfloor \approx 1.67\times10^{6}$ picks
and take $1440\times T \approx 3.07\times10^{8}$ trading steps.
\end{enumerate}

\section{Results}

\begin{center}
\renewcommand{\arraystretch}{1.2}
\begin{tabular}{@{}l r r@{}}
\toprule
Re-pick every $\Delta$ & $F(\Delta)$ & buys \\
\midrule
1 minute & \$4{,}376.71 & 16 \\
2 minutes & \$4{,}733.99 & 16 \\
20--39 minutes (all the same) & \$4{,}733.99 & 16 \\
60 minutes & \$4{,}733.99 & 16 \\
120 minutes & \$4{,}912.63 & 16 \\
424 minutes (highest, a lucky spike) & \$5{,}408.79 & 15 \\
1{,}440 minutes (one day) & \$3{,}732.82 & 14 \\
\midrule
Just holding & \$2{,}829.82 & -- \\
\bottomrule
\end{tabular}
\end{center}
999 of the 1,440 intervals beat holding. Intervals of 1 to 120 minutes average \$4{,}613.

\section{What the math does not cover}

\begin{itemize}[leftmargin=*]
\item One history. All numbers come from one five-month path of PEARL; a different path can rank
the settings differently.
\item Most of the machine's lead came from staying in USDC during the July--August fall.
From 1 Aug onward, holding grew more.
\item Fills are assumed at the next minute's close with a fixed $0.7\%$ cost. Real fills depend
on the order book at that moment.
\item The rule reads PEARL-USDT; the practice bots trade at the PEARL-USDC order book.
\end{itemize}

\appendix
\section{Other layers built along the way}

These were tested earlier and are not part of the machine above.

\paragraph{Live monitor.} Every second: mid price $m_t=(\text{bid}_t+\text{ask}_t)/2$, spread
$100(\text{ask}_t-\text{bid}_t)/m_t$, order-book imbalance
$(Q^{\text{bid}}_t-Q^{\text{ask}}_t)/(Q^{\text{bid}}_t+Q^{\text{ask}}_t)$ over the top 5 levels, and
UP/DOWN when the 10-reading average moves away from the 60-reading average by more than the spread.

\paragraph{Trend rule.} Hold PEARL while $\overline{c}^{(24\text{h})}_t > \overline{c}^{(168\text{h})}_t$
(simple moving averages of hourly closes).

\paragraph{Swing rule.} Buy when $c_t \ge \ell_t(1+r)$ where $\ell_t$ is the lowest price since the last
sell; sell when $c_t \le u_t(1-d)$ where $u_t$ is the highest price since the buy.

\paragraph{Floor rule.} After buying at $B$ with highest price $u_t$ since then: if
$u_t \ge 1.016B$, sell when $c_t \le \max(1.015B,\ u_t(1-\text{trail}))$; otherwise sell when
$c_t\le u_t(1-\text{stop})$.

\paragraph{Spike rule.} Over the last $W$ seconds with low $\ell$ and high $u$: a drop spike is
$p\le u(1-\text{drop\_spike})$ (buy, sell at $+$bounce or $-$bounce\_stop); a rise spike is
$p\ge \ell(1+\text{rise\_spike})$ (sell if holding, block buys 10 minutes).

\paragraph{Odds of a move continuing.} For a drop of $X\%$ below the latest top $u$: the chance
the price reaches $u(1-X)(1-0.016)$ before it gets back above $u$. For a rise of $X\%$ above the
latest bottom: the mirror image.

\paragraph{Why a 60/40 edge was not enough.} A trade wins $M$ with chance $p$ and loses $M$ otherwise:
expected gain $(2p-1)M$. It pays only if
\[
(2p-1)\,M \ \ge\ 1.4\% ,
\]
so $p=0.6$ needs moves of $M\ge 7\%$, and $p=0.8$ needs $M\ge 2.33\%$.

\end{document}
